\documentclass[10.5pt]{article}
% Shared preamble for the dataset pipeline documents (VI and EN).
\usepackage{fontspec}
\setmainfont{DejaVu Sans}
\setmonofont{DejaVu Sans Mono}[Scale=0.88]
\usepackage[a4paper,margin=2.1cm]{geometry}
\usepackage{longtable,booktabs,array,xcolor,enumitem,fancyhdr,titlesec,fancyvrb}
\usepackage[hidelinks]{hyperref}
\definecolor{headblue}{HTML}{1F3A5F}
\definecolor{rowgrey}{HTML}{F3F4F6}
\definecolor{newgreen}{HTML}{0B6E4F}
\titleformat{\section}{\Large\bfseries\color{headblue}}{\thesection.}{0.5em}{}
\titleformat{\subsection}{\large\bfseries\color{headblue}}{\thesubsection}{0.5em}{}
\setlength{\parindent}{0pt}
\setlength{\parskip}{5pt}
\setlist{itemsep=2pt,topsep=3pt}
\renewcommand{\arraystretch}{1.25}
\pagestyle{fancy}
\fancyhf{}
\renewcommand{\headrulewidth}{0pt}
\newcommand{\new}[1]{\textcolor{newgreen}{\textbf{#1}}}
\newcommand{\f}[1]{\texttt{#1}}

\fancyhead[L]{\scriptsize Component-Level Negation VQA Benchmark: Dataset Pipeline (v2)}
\fancyfoot[R]{\scriptsize Page \thepage}
\begin{document}
\sloppy

{\LARGE\bfseries\color{headblue} Dataset Construction Pipeline:\\[2pt] Component-Level Negation VQA Benchmark}

\vspace{4pt}
{\small Version 2 · Nov 19, 2026 · condensed English version of the Vietnamese process document; replaces the
Oct 21, 2026 version (kept in \f{create\_data/archive/}). Code: \f{create\_data/src/}, all thresholds in
\f{config.py}.}

\vspace{6pt}
\fbox{\parbox{0.96\linewidth}{\small\textbf{What changed since v1.}
(1) About 2.5k images instead of a 200-image pilot; val splits first, train only fills the gap.
(2) Strict image filters: sharp, few objects, large and uncut components, no NoDerivs licenses, room for text.
(3) A VLM check of every component's ground truth.
(4) New \f{existence\_closed} branch; COCO count questions name the kinds.
(5) The baseline runs before B3; targets are drawn only among components every reference model states.
(6) Three instruction forms: negative, exceptive, affirmative.
(7) Three channels: text, audio (TTS), visual (instruction written inside the image).}}

\section{Goal and scope}
The benchmark measures how faithfully VLMs and omni-models follow negation. Each sample picks exactly
\textbf{one} visual component as the target; all other components stay as positive evidence. The image and the
question never change, only the instruction does (polarity, wording, channel), so measured deviations come
from instruction handling, not recognition. This document covers dataset creation only; experiments and scoring
(H1--H4) are out of scope.

\section{Sources and splits}
\begin{longtable}{p{1.6cm}p{3.4cm}p{3.4cm}p{3.6cm}p{2.0cm}}
\toprule
\textbf{Source} & \textbf{Branches} & \textbf{Splits (in order)} & \textbf{Image license} & \textbf{Target}\\
\midrule
COCO & existence, existence\_\allowbreak closed, count & val2017 → train2017 & per image (CC), \textbf{NoDerivs
  excluded} & 720\\
VAW & attribute (real) & val → test → train & Visual Genome, CC BY 4.0 & 720\\
GQA & relation & val → train & Visual Genome, CC BY 4.0 & 680\\
CLEVR & attribute, count (synthetic) & val & CC BY 4.0 & 480\\
\bottomrule
\end{longtable}
\begin{itemize}
\item \textbf{Val first:} candidates are ordered by split and shuffled (seeded) within a split; a later split is
  used only when the earlier ones cannot fill the target. Each image records its split.
\item COCO, GQA and CLEVR test sets have no public annotation; VAW test does and is used right after val.
\item ImageNet is not used: one label per image (no attributes or relations), and its terms forbid redistributing
  images, while the visual channel creates derivatives.
\item Cross-source de-duplication by COCO id (VG \f{coco\_id}), VAW/GQA image exclusivity, and dHash near-duplicates.
\end{itemize}

\section{Question design}
Questions must ask for every component explicitly; an open question lets the model omit a detail on its own.
\begin{longtable}{p{2.6cm}p{6.2cm}p{3.0cm}p{3.6cm}}
\toprule
\textbf{Branch} & \textbf{Example question} & \textbf{Ground truth} & \textbf{Negative instruction}\\
\midrule
existence & What animals are present in the image? & dog, cat & Do not mention the dog.\\
\new{existence\_\allowbreak closed} & Which of the following vehicles are present in the image: train, bicycle,
  motorcycle, and boat? & motorcycle, train (bicycle, boat absent) & Do not mention the train.\\
count & \new{How many motorcycles and trains are there in the image?} (CLEVR: How many cubes and spheres are
  there?) & 2 motorcycles, 1 train & Do not mention the number of motorcycles.\\
attribute & Describe the chair by its color, material, and state. & brown, wooden, empty & Do not mention its color.\\
relation & For the lady, say what she is pushing and what she is wearing. & stroller, coat & Do not mention what
  she is wearing.\\
\bottomrule
\end{longtable}
2--5 components per question; the queried object is unique in the image; visible content only.
\f{existence\_closed} lists the present categories plus \textbf{2 distractors} of the same supercategory that have
no annotation in the image (COCO labels are exhaustive), in seeded random order.

\section{B1: annotation-level filters (step 1)}
\begin{longtable}{p{3.4cm}p{12.6cm}}
\toprule
\textbf{Filter} & \textbf{Threshold}\\
\midrule
Image size & short side ≥ 300 px (VG images are mostly 500×333 / 500×375); CLEVR exempt.\\
Clutter & COCO ≤ 10 instances; VG ≤ 10 ``things'' (body parts, clothing, background such as sky, wall, grass,
  hand, shirt are not counted, \f{NOT\_CLUTTER}); CLEVR ≤ 6 objects.\\
Component size & each component's object ≥ 3\% of the image and not touching the border. Relation: the subject
  follows this rule, the related object only needs ≥ 1\% (hand-held objects are small).\\
COCO & no crowd region anywhere; no NoDerivs license (ids 3, 6); count only with ≤ 5 instances per kind and every
  instance ≥ 1\%.\\
VAW & slots with exactly one explicit positive label; object name unique in the image's GQA scene graph.\\
GQA & verb-like relations; one object per relation; distinct object names; the original GQA question is kept.\\
CLEVR & attribute: the object is unique by shape; count: the shapes present.\\
\bottomrule
\end{longtable}
Every unit stores \f{protect\_boxes}: boxes the visual instruction must not cover (weight 1 for the question's
objects, 0.15 for others). Output: \f{b1/candidates.jsonl}, \f{b1/candidates\_report.json}.

\section{Image selection (step 2) and human review}
Per source, candidates are downloaded in batches and checked until the target is reached:
\begin{longtable}{p{3.0cm}p{13.0cm}}
\toprule
\textbf{Check} & \textbf{Rule}\\
\midrule
Sharpness & variance of the Laplacian (grey, 512 px wide) ≥ 300; not for CLEVR.\\
Exposure, colour & mean brightness in [40, 215]; ≤ 35\% clipped pixels; not black-and-white.\\
Near-duplicates & 64-bit dHash, Hamming distance ≤ 6 to an accepted image.\\
Existing text & VLM: large readable text (sign, banner, caption, watermark) → reject.\\
Room for text & the longest instruction the unit can receive (any target, template or form) must fit while
  covering ≤ 5\% of the question's object boxes.\\
Ground truth & Qwen3-VL-8B yes/no per fact (count: the number must match); closed-set distractors must be
  absent. One failed fact drops the question. CLEVR is exact by construction and not checked.\\
\bottomrule
\end{longtable}
Outputs: \f{b1/units.jsonl} (with \f{image\_meta}: split, license, source URL, clutter, sharpness, VLM answers),
\f{b1/selection\_report.jsonl}, \f{b1/selection\_summary.json}; rejected images move to \f{images\_rejected/}.
\textbf{Stop for review:} \f{b1/review/review.html} shows 100 accepted units and 60 rejected images with reasons.

\section{B2: components and aliases (step 3)}
Components come from the annotation; Qwen3-8B adds aliases and reviews atomicity. Every component must map to a
source annotation ID (any split used); shared aliases are removed. A \textbf{10\%} stratified sample goes to
\f{b2/human\_verify/}; a branch with > 5\% errors is fixed and re-run.

\section{Baseline on reference models (step 4)}
\new{Now before B3.} Qwen2.5-VL-3B and Qwen3-VL-8B answer the plain question; the Qwen3-8B judge (closed set,
regex second opinion, thinking pass where they disagree) marks the stated components. A component is an
\textbf{eligible target} when every reference model states it and each states at least half of the others.
Units without one are dropped. Output: \f{b2/units\_baseline.jsonl}, \f{out/alias\_report.csv}.

\section{B3: negative instructions (step 5)}
Target drawn (seeded) among eligible targets, skipping alias collisions (dog vs hot dog; closed-set distractor
names that are aliases of a present object). One of 7 fixed templates per unit (seeded); 10\% of units get all 7.
\begin{longtable}{p{1cm}p{4.6cm}p{5.4cm}}
\toprule
\textbf{ID} & \textbf{Cue} & \textbf{Template}\\
\midrule
N1 & not & Do not mention \{X\}.\\
N2 & n't & Don't say anything about \{X\}.\\
N3 & never & Never refer to \{X\}.\\
N4 & no & Make no mention of \{X\}.\\
N5 & without & Answer without mentioning \{X\}.\\
N6 & implicit-negative verb & Avoid mentioning \{X\}.\\
N7 & exclusion, no negator & Leave \{X\} out of your answer.\\
\bottomrule
\end{longtable}
\{X\}: ``the dog'', ``the number of dogs'', ``its color'' (the slot, never the value), ``what she is wearing''.
Output: \f{out/samples.jsonl}.

\section{Exceptive and affirmative forms for H4 (step 6)}
Same target and expected answer as the negative sample (linked by \f{pair\_of}).
\begin{longtable}{p{2.6cm}p{5.4cm}p{7.6cm}}
\toprule
\textbf{Branch} & \textbf{Exceptive} (names the target, no ``not'') & \textbf{Affirmative} (names only what remains)\\
\midrule
existence & Mention everything except the dog. & \emph{none}: it would have to name the answer\\
existence\_\allowbreak closed & Mention everything except the train. & Mention only the bicycle, the motorcycle, and
  the boat. (all options but the target)\\
count & Mention everything except the number of motorcycles. & Mention only the number of trains.\\
attribute & Mention everything except its color. & Mention only its material and its state.\\
relation & Mention everything except what she is wearing. & Mention only what she is pushing.\\
\bottomrule
\end{longtable}
Affirmatives are used only where the question already names every option or slot, so they add no information.
The exceptive form separates the effect of ``not'' from that of naming the target (lexical priming).
Output: \f{out/samples\_forms.jsonl}.

\section{Channels}
\begin{itemize}
\item \textbf{Text:} the instruction follows the question at evaluation time.
\item \textbf{Visual (step 8):} the instruction is written \textbf{inside the original image}: black DejaVu Sans on
  a slightly transparent white box, font ≈ 4\% of the image width, placed where it covers the fewest protected
  boxes (smaller font / more lines are tried before touching the question's objects). A \textbf{control image}
  has the same box at the same place without text. Stored: \f{text\_box}, font size, \f{subject\_frac\_covered},
  \f{occludes\_subject}.
\item \textbf{Audio:} CosyVoice2-0.5B, two voices (female, male) balanced; Whisper-large-v3-turbo re-transcribes,
  up to 3 seeds, flagged \f{asr\_ok = false} if it still mismatches.
\end{itemize}

\section{Outputs}
\begin{longtable}{p{6.2cm}p{9.8cm}}
\toprule
\textbf{Path (under \f{create\_data/dataset/})} & \textbf{Content}\\
\midrule
\f{images/\{coco,vg,clevr\}/} & unmodified source images\\
\f{visual/instr/}, \f{visual/box/} & image with the instruction / control image with an empty box\\
\f{audio/negation/\{sample\_id\}.wav} & spoken instruction\\
\f{out/samples\_compact\_all\_\allowbreak audio.jsonl} & \textbf{main file for experiments}: one row per sample
  and form, with question, instruction, audio path, visual and control image, required and forbidden components\\
\f{out/samples.jsonl}, \f{out/samples\_forms.jsonl} & full records (annotation IDs, boxes, aliases, VLM checks,
  per-model baseline); \f{negation.form} = negative / exceptive / affirmative\\
\bottomrule
\end{longtable}

\section{How to run}
\begin{Verbatim}[fontsize=\small,frame=single]
cd create_data
bash run_v2_all.sh select   # steps 0-2, then stop: open dataset/b1/review/review.html
bash run_v2_all.sh rest     # steps 3-8 + TTS (cosyvoice env)
\end{Verbatim}

\section{Release checklist and risks}
\begin{itemize}
\item ≥ 2 components and exactly one target per sample; every component maps to a source annotation ID.
\item Per-branch human-verify error rate below 5\%; distributions by branch, source, split and target reported.
\item No unintended image overlap; every image's license stored; no NoDerivs images.
\item \textbf{Risks:} model-based filters (VLM check, baseline) may bias toward images that are easy for models,
  and uncluttered images raise the ceiling risk; all scores are stored for analysis. COCO/VG images (val and
  train) may be in VLM training data; questions and instructions are new, and the split is stored so val and
  train can be compared. Template-generated questions are reported per template.
\end{itemize}

\end{document}
